% !TeX root = ../../tesis.tex
%%%%%%%%%%%%%%%%%%%%%%%%%%%%%%%%%%%%%%%%%%%%%%%%%%%%%%%%%%%%%%%%%%%
% SECCIÓN 5.3.1 — VALIDACIÓN DEL MOTOR DE IMPEDANCIA
%%%%%%%%%%%%%%%%%%%%%%%%%%%%%%%%%%%%%%%%%%%%%%%%%%%%%%%%%%%%%%%%%%%
\subsection{Validación del Motor de Impedancia Temporal}
\label{subsec:validacion-impedancia}

Para verificar la consistencia del motor de impedancia de \texttt{VFTModel},
se calculó la velocidad efectiva implícita de cada sistema tras aplicar el
Coeficiente de Fricción Vial ($CF$) y se contrastó con rangos operativos de
referencia por sistema. La velocidad efectiva se obtiene a partir de la
Ecuación~\eqref{eq:velocidad-efectiva}:

\begin{equation}
    V_{\text{efectiva}} = \frac{V_{\text{base}}}{CF}
    \label{eq:velocidad-efectiva}
\end{equation}

\noindent donde:
\begin{itemize}
    \item $V_{\text{base}}$: velocidad promedio de operación del sistema, en
          km/h, registrada por el \texttt{Apímetro} para cada ruta.
    \item $CF$: coeficiente de fricción según el tipo de derecho de vía
          (véase §\ref{subsec:fase2-cf}).
\end{itemize}

En el cuadro, la velocidad efectiva de cada sistema es el promedio de la
relación distancia/tiempo de sus segmentos; en los sistemas con un solo tipo de
derecho de vía ese promedio coincide con la
Ecuación~\eqref{eq:velocidad-efectiva}. El Cuadro~\ref{tab:validacion-impedancia}
sintetiza los resultados de la corrida ejecutada sobre
\textbf{25{,}025 segmentos de servicio real}. El Trolebús se presenta en dos
filas porque su Línea~10, elevada, opera en vía confinada.

\begin{table}[H]
    \centering
    \caption[Validación de velocidades efectivas por sistema de transporte]{
    Contraste entre la velocidad implícita post-fricción calculada por
    \texttt{VFTModel} y los rangos operativos nominales de cada sistema.
    Corrida: 2026-10-04.}
    \label{tab:validacion-impedancia}
    \begin{tabular}{|l|r|r|c|c|}
        \hline
        \textbf{Sistema} & \textbf{$V_{\text{base}}$ (km/h)} &
        \textbf{$V$ efectiva (km/h)} &
        \textbf{Rango esperado} & \textbf{Estado} \\
        \hline
        Tren Interurbano       & 70.0 & 70.0 & ---           & Sin rango \\
        Tren Suburbano         & 65.0 & 65.0 & {[}55--75{]}  & \checkmark\ OK \\
        Metro (STC)            & 36.0 & 36.0 & {[}30--42{]}  & \checkmark\ OK \\
        Tren Ligero            & 22.0 & 22.0 & {[}18--26{]}  & \checkmark\ OK \\
        Trolebús (Línea 10)    & 25.0 & 21.7 & {[}14--22{]}  & \checkmark\ OK \\
        Cablebús               & 20.0 & 20.0 & {[}16--24{]}  & \checkmark\ OK \\
        Mexicable              & 20.0 & 20.0 & {[}16--24{]}  & \checkmark\ OK \\
        Metrobús               & 16.3 & 16.3 & {[}13--20{]}  & \checkmark\ OK \\
        Mexibús                & 16.3 & 14.2 & {[}13--20{]}  & \checkmark\ OK \\
        Trolebús (compartido)  & 18.0 & 13.0 & {[}14--22{]}  & $\approx$ Límite \\
        Pumabús                & 14.0 & 10.1 & {[}8--14{]}   & \checkmark\ OK \\
        \textsc{rtp}           & 11.0 &  8.0 & {[}10--18{]}  & Bajo rango \\
        Corredores Conc.       & 11.0 &  8.0 & {[}8--14{]}   & $\approx$ Límite \\
        \hline
    \end{tabular}
    \fuentefigura{Elaboración propia con \texttt{VFTModel}
    \citep{cabrera2026VFTModel}. Fuentes de referencia: STC Metro \cdmx{},
    Metrobús \cdmx{}, \textsc{rtp} \cdmx{}, \textsc{ste} Trolebús.}
\end{table}

\begin{figure}[H]
    \centering
    \includegraphics[width=\textwidth]{Figures/Cap5/fig01_impedancia.png}
    \caption[Caracterización de la impedancia temporal por sistema]{(a)~Velocidad
    efectiva post-fricción por sistema, coloreada según el tipo de derecho de
    vía. (b)~Distribución de tiempos de viaje por segmento interestación
    (segmentos menores a 5~min, sin valores atípicos); el Tren Interurbano no
    aparece en este panel porque sus cuatro segmentos superan ese umbral.
    Corrida: 2026-10-04.}
    \label{fig:impedancia-sistemas}
    \fuentefigura{Elaboración propia con \texttt{VFTModel}
    \citep{cabrera2026VFTModel}.}
\end{figure}

Los resultados son coherentes con la metodología adoptada en los sistemas de vía
exclusiva y confinada, que quedan dentro de sus rangos. Las desviaciones se
concentran en el transporte de superficie. El \textbf{Trolebús} en vía
compartida alcanza $\approx$13.0~km/h con $CF = 1.380$, alrededor de 1~km/h (7\,
\%) por debajo del límite inferior de su rango (14--22~km/h); el tipo de derecho
de vía \termino{compartido} no captura la preferencia operativa parcial que el
trolebús mantiene en ciertas vialidades de la \cdmx{}. La \textsc{rtp} y los
Corredores Concesionados parten de una velocidad base de 11~km/h y, con el mismo
$CF$, llegan a 8.0~km/h: los Corredores quedan en el límite inferior de su rango
y la \textsc{rtp} unos 2~km/h por debajo del suyo. Para el análisis estructural
de la red estas velocidades representan un escenario conservador del transporte
de superficie; contrastarlas con velocidades observadas por ruta queda como
trabajo futuro.
